% ID: MAT-TRI-EQG-003
% Titolo: Equazioni goniometriche con fattorizzazione (3)
% Disciplina: Matematica
% Area: Trigonometria
% Argomento: Equazioni goniometriche
% Classe: Quarta secondaria di secondo grado
% Difficolta: 4
% Tipo: problema_numerico
% Tag: matematica, trigonometria, equazioni_goniometriche
% Fonte: verifica_goniometria_avanzata_anonimizzata
% Provenienza: materiale fornito dal docente; testo anonimizzato; soluzione redatta per Exergo

\begin{esercizio}
Risolvi le seguenti equazioni: 

\begin{enumerate}
\item $\displaystyle
\sin^2x+(\sqrt{3}-1)\sin x\cos x-\sqrt{3}\cos^2x=0;$

\item $\displaystyle
\dfrac{\sin 2x}{1-\sin x}=2\cos x.$
\end{enumerate}
Salvo diversa indicazione, lavora in $\mathbb R$; angoli in radianti.
\end{esercizio}

\begin{soluzione}
a) Per $\cos x=0$ il primo membro vale 1. Altrimenti, dividendo per $\cos^2x$,
$(\tan x-1)(\tan x+\sqrt3)=0$: $x=\pi/4+k\pi$ oppure $x=-\pi/3+k\pi$.
b) Il dominio esclude $\sin x=1$. Moltiplicando per $1-\sin x$ si ricava
$2\cos x(2\sin x-1)=0$. Restano $x=3\pi/2+2k\pi$, $x=\pi/6+2k\pi$, $x=5\pi/6+2k\pi$. In tutte le famiglie di soluzioni $k\in\mathbb{Z}$.
\end{soluzione}
